% ============================================================================
%  OPEN ITEMS REGISTER -- Cubli Technical Design Document
% ----------------------------------------------------------------------------
%  Standalone working document. NOT input by main.tex -- it is a checklist for
%  the team and a briefing sheet to hand to an assistant, not part of the TDD.
%
%  Compile on its own:  pdflatex OPEN_ITEMS.tex
%
%  Scope: every \textcolor{red}{TBD} in sections/*.tex, every empty section,
%  every comment-only stub and every missing figure, as of the scan date below.
%
%  RESTRUCTURE (second/final draft): the document was reorganised from
%  Introduction / Project Organisation / System Analysis / Dynamic Modelling &
%  Control / Preliminary Hardware Design / Future Work into Introduction /
%  Project Organisation / Preliminary Sizing / Methodology / CAD Design /
%  Electronics / Control Algorithms & Approach / Conclusions & Next Steps.
%  Section~\ref{sec:filemap} maps old content to its new home. The two
%  previously BLOCKER content gaps -- Appendix D (control algorithms) at 0
%  bytes, and Section 3 (System Analysis) three-quarters empty -- are both
%  CLOSED structurally: real prose and tables exist at both locations now.
%  What remains open at those locations is tracked below as ordinary TBD
%  markers (mostly gains, timing and physical pin numbers), not as missing
%  sections. Uncommenting the electrical schematic-sheet placeholders and
%  adding a per-pin logic table to each also added a new block of TBD figures
%  (the physical pin numbers, pending the board layout) that did not exist at
%  the previous scan.
% ============================================================================
\documentclass[11pt,a4paper]{article}

\usepackage[utf8]{inputenc}
\usepackage[T1]{fontenc}
\usepackage[a4paper,margin=2.2cm]{geometry}
\usepackage{amsmath,amssymb}
\usepackage{booktabs}
\usepackage{longtable}
\usepackage{array}
\usepackage{ragged2e}
\usepackage{xcolor}
\usepackage{siunitx}
\usepackage{fancyhdr}
\usepackage{lastpage}
\usepackage{hyperref}
\usepackage{caption}
\usepackage{parskip}

\hypersetup{colorlinks=true, linkcolor=blue, urlcolor=blue,
            pdftitle={Cubli TDD - Open Items Register}}

\DeclareSIUnit{\rpm}{rpm}
\newcolumntype{L}[1]{>{\RaggedRight\arraybackslash}p{#1}}
\renewcommand{\arraystretch}{1.15}
\captionsetup{font=small, labelfont=bf, skip=6pt}

% Status shorthands, so the priority language is consistent across the tables.
\definecolor{blockcol}{HTML}{C00000}
\definecolor{readycol}{HTML}{70AD47}
\definecolor{waitcol}{HTML}{BF8F00}
\definecolor{closedcol}{HTML}{2E75B6}
\newcommand{\blocker}{\textcolor{blockcol}{\textbf{BLOCKER}}}
\newcommand{\ready}{\textcolor{readycol}{\textbf{READY}}}
\newcommand{\waiting}{\textcolor{waitcol}{\textbf{WAITING}}}
\newcommand{\closed}{\textcolor{closedcol}{\textbf{CLOSED}}}

\pagestyle{fancy}
\fancyhf{}
\setlength{\headheight}{14pt}
\fancyhead[L]{\textbf{Open Items Register}}
\fancyhead[R]{Cubli TDD}
\fancyfoot[R]{Page \thepage\ of \pageref{LastPage}}
\renewcommand{\headrulewidth}{0.4pt}
\renewcommand{\footrulewidth}{0.4pt}

\title{\textbf{Open Items Register}\\[0.4em]\large Technical Design Document --- Cubli}
\author{Armonia --- Space Challenges 2026}
\date{Scan of \today}

\begin{document}
\maketitle
\thispagestyle{fancy}

\section*{How to use this document}
\addcontentsline{toc}{section}{How to use this document}

This register lists every unresolved item in the TDD: the \textcolor{red}{TBD}
markers, the sections that are empty, the subsections that exist only as
comments, and the figures that are referenced but not drawn. It is a working
checklist, not part of the submitted document.

Three status labels are used throughout:

\begin{itemize}
  \item \ready{} --- decidable now. Nothing upstream is missing; the item is
        waiting on a decision or a datasheet lookup, not on data.
  \item \waiting{} --- blocked on an upstream item, named in the row.
  \item \blocker{} --- blocks a large number of downstream items. Close these
        first.
  \item \closed{} --- resolved since the original scan. Retained in the register
        so the closure is auditable rather than silently deleted.
\end{itemize}

Section~\ref{sec:filemap} gives the section-file map for the current draft.
Section~\ref{sec:chain} states the dependency chain, which is what makes the
ordering non-arbitrary: roughly a quarter of the register closes automatically
once the packaging study closes.

\section{Section-file map (current draft)}
\label{sec:filemap}

\begin{table}[h]
\centering
\small
\begin{tabular}{c L{4.4cm} L{7.0cm}}
\toprule
\# & Section & File \\
\midrule
1 & Introduction (incl.\ RW-vs-CMG trade) & \texttt{01\_introduction.tex} \\
2 & Project Organisation (incl.\ validation stage-gate table) & \texttt{02\_project\_organization*.tex} \\
3 & Preliminary Sizing & \texttt{03\_preliminary\_sizing.tex} \\
4 & Methodology (incl.\ verification/checks) & \texttt{04\_methodology.tex} \\
5 & CAD Design (incl.\ CAD verification) & \texttt{05\_cad\_design.tex} \\
6 & Electronics & \texttt{06\_electronics.tex} \\
7 & Control Algorithms \& Approach & \texttt{07\_control\_algorithms.tex} \\
8 & Conclusions \& Next Steps (incl.\ validation results, future work) & \texttt{08\_conclusions.tex} \\
A & Bill of Materials & \texttt{09\_bom.tex} \\
B & CAD \& Construction Detail & \texttt{10\_cad\_detail.tex} \\
C & Electrical System Detail (schematic sheets, now uncommented, plus per-pin tables) & \texttt{11\_electrical\_detail.tex} \\
D & Control Algorithm Detail (state machine, pseudocode, gains, timing, faults) & \texttt{12\_control\_algorithm\_detail.tex} \\
\bottomrule
\end{tabular}
\end{table}

The former \texttt{03\_analysis.tex} (RW-vs-CMG trade, requirements
traceability) and \texttt{06\_testing.tex} (excluded from the build) are
retired; their content is redistributed into the files above rather than
existing as standalone sections. Section labels most items are still
addressed by (\texttt{sec:massbudget}, \texttt{sec:jumpup\_target},
\texttt{sec:rw\_sizing}, \texttt{sec:control2d}, \texttt{sec:test\_feasibility}
$\to$ \texttt{sec:sizing\_verification}, \texttt{sec:test\_electrical} $\to$
\texttt{sec:elec\_checks}, \texttt{sec:test\_cad} $\to$
\texttt{sec:cad\_verification}, \texttt{sec:test\_sim} $\to$
\texttt{sec:sim\_environment}) are unchanged or carried forward under a new
name; the label, not the file, is the stable identifier to search on.

\section{Dependency chain}
\label{sec:chain}

Most items are not independent. The governing chain is

\begin{center}
$V_{\min} \;\rightarrow\; L \;\rightarrow\; M \;\rightarrow\; h_w \;\rightarrow\;
I_{w,\text{target}} \;\rightarrow\;$ wheel design, feasibility gate, test thresholds
\end{center}

so the packaging study (item \textbf{A4}, Table~\ref{tab:cad_params} in
\texttt{10\_cad\_detail.tex}) is the true head of the register. The
parametric calculator \texttt{analysis/SIZING.py} implements Stages~1--6 of the
closure procedure (Section~\ref{sec:massbudget}, now in
\texttt{03\_preliminary\_sizing.tex}) and reports a converged fixed point at
$M = \qty{1.7716}{\kilo\gram}$, reproduced by its own Stage~6 roll-up to within
\qty{0.1}{\gram}. The figures below are \textbf{transcribed into the TDD}
(A1, A2 and A8 are closed). They remain provisional because roughly
\qty{34}{\percent} of the mass total is still structural allowance rather than
measured mass.

\begin{table}[h]
\centering
\caption[Values currently produced by the sizing calculator]{Converged values
from \texttt{analysis/SIZING.py}, as transcribed into the TDD. The wheel
diameter is pinned by the packaging envelope
(\texttt{D\_w\_FIXED\_mm}\,=\,120), not chosen by the mass-minimising sweep.}
\label{tab:sizing_current}
\small
\begin{tabular}{l L{5.0cm} c c}
\toprule
Symbol & Quantity & Converged value & Status \\
\midrule
$M$                   & Cube mass (Stage 1 assumption) & \qty{1.7716}{\kilo\gram} & Fixed point \\
$M_{\text{est}}$      & Cube mass (Stage 6 roll-up)    & \qty{1.7716}{\kilo\gram} & $\Delta \approx 0$; 34\% allowance \\
$L$                   & Cube edge length               & \qty{150}{\milli\metre} & Fixed \\
$\eta$                & Energy margin                  & 1.05                    & Chosen \\
$\tau_b$              & Brake torque per wheel         & \qty{5}{\newton\metre}  & Assumed \\
$\tau_g = MgL/2$      & Gravitational tip-over floor   & \qty{1.303}{\newton\metre} & Derived \\
$\beta$               & Brake-amplification factor     & 1.353                   & Derived \\
$\omega_{\max}$       & Max wheel speed                & \qty{6000}{\rpm}        & Agreed with TDD \\
$h_w$                 & Required angular momentum      & \qty{0.2588}{\kilo\gram\metre\squared\per\second} & Derived \\
$I_{w,\text{target}}$ & Target wheel inertia (corner)  & \qty{4.119e-4}{\kilogram\metre\squared} & Derived \\
$I_{w,\text{target}}$ & Target wheel inertia (edge)    & \qty{3.737e-4}{\kilogram\metre\squared} & Not critical \\
$D_w$                 & Wheel outer diameter           & \qty{120}{\milli\metre} & Envelope constraint \\
$N$                   & Ballast stations per wheel     & 15 of $N_{\max}=39$     & Derived \\
$m_w$                 & Wheel mass (each)              & \qty{172.5}{\gram}      & Derived \\
$I_w$                 & Wheel inertia achieved         & \qty{4.599e-4}{\kilogram\metre\squared} & \qty{111.7}{\percent} of target \\
\bottomrule
\end{tabular}
\end{table}

\paragraph{What remains open on the sizing chain.} The residual risk is not a
conflict but a margin asymmetry: the wheel sits \qty{11.7}{\percent} above
target while one station is worth \qty{5.9}{\percent}, so downward ballast trim
is nearly exhausted ($-1$ station) while upward trim is ample ($+24$). A
structure that comes in \emph{lighter} than the \qty{610}{\gram} allowance is
therefore the awkward case, and $\omega_{\max}$ is the relief variable. This is
unchanged by the restructure --- the numbers moved files, not values.

\clearpage
\section{TBD items, by owning table}
\label{sec:tbd}

\subsection{A1 --- Jump-up sizing budget \quad \closed}

\texttt{tab:jumpup\_budget}, \texttt{03\_preliminary\_sizing.tex}
(\texttt{sec:jumpup\_target}). All 6 values populated from the converged run
of Table~\ref{tab:sizing_current}; provisional in the sense that $M$ carries a
\qty{34}{\percent} structural allowance.

\subsection{A2 --- Required inertia by contact case \quad \closed}

\texttt{tab:inertia\_cases}, \texttt{03\_preliminary\_sizing.tex}
(\texttt{sec:massbudget}). Both contact cases populated: edge
$I_{w,\text{target}} = \qty{3.737e-4}{\kilogram\metre\squared}$; corner
\qty{4.119e-4}{}. The corner case demands \qty{9.3}{\percent} more inertia.

\subsection{A3 --- Rail current budget \quad \ready}

\texttt{tab:power\_budget}, \texttt{06\_electronics.tex}
(\texttt{sec:power\_arch}). \textbf{14 TBD}: typical and worst case for each of
five loads plus the logic-rail total and the driver spin-up figure. This table
was previously wrapped in a LaTeX \texttt{comment} environment and did not
render at all, despite being referenced from Appendix~C; it is now a live
table with the same TBD content. Typical values are datasheet lookups and can
be filled today; worst case is the \emph{coincident} case (telemetry burst,
servo stall and full sensor load together), not the sum of typicals.

\subsection{A4 --- Driving CAD parameters \quad \blocker}

\texttt{tab:cad\_params}, \texttt{10\_cad\_detail.tex}. \textbf{6 TBD,
unchanged}: $d_m$, $h_m$, $c_w$, $c_o$, $V_{\min}$, $r_c$. These need the
packaging study, not the calculator. $V_{\min}$ is step~2 of the closure
procedure and therefore the head of the whole dependency chain: it fixes $L$,
which fixes $M$, which fixes $I_{w,\text{target}}$. $c_w$ and $c_o$ are
additionally consumed by the interference checks of
\texttt{sec:cad\_verification} in \texttt{05\_cad\_design.tex}.

\subsection{A5 --- CAD model tree and part list \quad \ready}

\texttt{tab:cad\_parts}, \texttt{10\_cad\_detail.tex}. \textbf{34 TBD}:
fabrication route and status for every part across the inner structure, wheel
assembly, outer frame and 1-DoF rig, unchanged from the previous scan. All are
decidable today and do not depend on the sizing.

\subsection{A6 --- Print and manufacturing settings \quad \ready}

\texttt{tab:print\_settings}, \texttt{10\_cad\_detail.tex}. \textbf{18 TBD}:
material, infill/walls and orientation for the six structural printed parts
(W1, W2, M1, M2, F1, F2). The calculator assumes PET-CF at
\qty{1290}{\kilo\gram\per\metre\cubed}, which should be the material entry for
W1 at minimum.

\subsection{A7 --- Mass properties verification \quad \waiting}

\texttt{tab:mass\_verification}, \texttt{10\_cad\_detail.tex}. \textbf{12
TBD}: CAD value, measured value and discrepancy for $M$, $m_w$, $I_w$ and CoM
offset. The CAD column can be filled once the model exists; the measured
column requires assembled hardware.

\subsection{A8 --- Pre-fabrication feasibility gate \quad \closed}

\texttt{tab:jumpup\_gate}, \texttt{03\_preliminary\_sizing.tex}
(\texttt{sec:sizing\_verification}, formerly \texttt{sec:test\_feasibility}).
All 8 values populated and the gate passes on every row: speed margin
\num{0.896}; inertia route disagreement \qty{+11.7}{\percent}; ballast trim
$-1/+24$ stations; retention at \qty{1.47}{\percent} of proof load.

\subsection{A9 --- 1-DoF plant identification \quad \waiting}

\texttt{tab:paramid\_methods}, \texttt{04\_methodology.tex}
(\texttt{sec:paramid}). \textbf{4 TBD}: torque constant $K_t$, body inertia
$J$, friction, loop latency. Methods and instruments are specified; only the
measured values and tolerances are pending. Requires the rig.

\subsection{A10 --- Control performance metrics \quad \waiting}

\texttt{tab:perf\_metrics}, \texttt{08\_conclusions.tex}
(\texttt{sec:test\_perf}). \textbf{5 TBD}: targets for recovery angle,
settling time, disturbance-rejection bandwidth, slew tracking error, momentum
drift. Control effort is already defined by reference to the driver limits.

\subsection{A11 --- Connector schedule \quad \ready}

\texttt{tab:connectors}, \texttt{11\_electrical\_detail.tex}. \textbf{11 TBD}:
housing and keying for J1--J11, unchanged. No two connectors carrying
different voltages may share a housing type, so a mis-mate is impossible by
construction; J1 is fixed by the BOM (XT60).

\subsection{A12 --- Wire schedule \quad \ready}

\texttt{tab:wire\_schedule}, \texttt{11\_electrical\_detail.tex}. \textbf{14
TBD}: gauge and length for H1--H7, unchanged. Gauge follows from A3; length
follows from the CAD envelope.

\subsection{A13 --- As-built electrical measurements \quad \waiting}

\texttt{tab:electrical\_measured}, \texttt{11\_electrical\_detail.tex}.
\textbf{7 TBD}: $K_m$, $I_w$ from spin-down, viscous and Coulomb friction,
command-to-encoder latency, achievable wheel speed on the bus, worst-case
logic rail load, runtime per charge. Requires hardware.

\subsection{A14 --- Per-sheet pin/net logic tables \quad \ready\ / \waiting (mixed)}

\texttt{tab:pins\_power\_entry} through \texttt{tab:pins\_servo},
\texttt{11\_electrical\_detail.tex}, Sheets~1--7. \textbf{New since the
restructure}: uncommenting the schematic-sheet placeholders and adding a
per-pin electrical/logic table to each (net name, direction, level, function)
surfaced \textbf{9 TBD} physical-pin-number entries (E-stop/fuse pin
positions, the LM2596 pin-out, the CAN termination component refs, the Teensy
CAN-FD peripheral pins). The net identity, direction and level in every row
are fixed by the architecture of \texttt{sec:electrical} and are \emph{not}
TBD; only the physical pin number is, and it is \ready{} in the sense that it
needs a board layout decision (\texttt{sec:layout}), not upstream data --- but
that decision itself has not been made, so it is listed \waiting{} on it.

\subsection{A15 --- Appendix D: gains, timing budget, guard conditions \quad \waiting}

\texttt{12\_control\_algorithm\_detail.tex}. \textbf{New since the
restructure}: Appendix D was 0 bytes at the previous scan (see
Section~\ref{sec:filemap} note) and is now written in full --- the mode-
transition guard conditions (\texttt{tab:state\_transitions}, 5 TBD), the
fixed gains and tuning parameters (\texttt{tab:gains}, 10 TBD: $Q$, $R$,
$K_d$ for both the edge and corner models, $\gamma$, $\alpha$,
$\theta_{\text{th}}$, the desaturation threshold), the per-cycle timing budget
(\texttt{tab:timing\_budget}, 12 TBD across five pipeline stages), and two
fault-response timeouts (\texttt{tab:fault\_responses}, 2 TBD). All of these
close only once the plant identification (A9) and the Simulink/Simscape
simulation work of \texttt{sec:sim\_environment} report measured parameters
and re-derived gains --- they are a consequence of A9 closing, not an
independent blocker.

\section{TBD in prose}

\begin{table}[h]
\centering
\small
\begin{tabular}{L{4.4cm} L{9.2cm}}
\toprule
Location & Pending \\
\midrule
\texttt{08\_conclusions.tex}, \texttt{sec:test\_edge} & Edge-balance numeric
targets beyond the G6 gate criterion. \\
\texttt{08\_conclusions.tex}, \texttt{sec:test\_slew} & Slew tracking
targets. \\
\texttt{07\_control\_algorithms.tex}, \texttt{sec:sim\_environment} &
Robustness-sweep perturbation magnitude --- set once the identification
tolerances of A9 are known. \\
\bottomrule
\end{tabular}
\end{table}

\section{Empty sections}

\textbf{None.} Both previously-empty content blocks are closed: Appendix D
(\texttt{12\_control\_algorithm\_detail.tex}) now carries the mode state
machine, estimator and controller pseudocode, gains table, timing budget and
fault-response table (residual TBDs tracked as A15); and the former Section~3
(RW-vs-CMG trade, requirements traceability) is folded into the Introduction
and the opening of Preliminary Sizing respectively, so there is no longer a
stub section for it.

Three headings remain deliberately empty, by standing decision rather than by
omission --- they are commented-out placeholders reserved for content not yet
written, per the placeholder convention used throughout the document
(numbered heading + label + one-line note, no filler prose):

\begin{table}[h]
\centering
\small
\begin{tabular}{L{4.6cm} L{4.0cm} L{5.0cm}}
\toprule
Section & File & State \\
\midrule
NMPC \& warm-start control   & \texttt{08\_conclusions.tex} & Heading only (commented) \\
Autonomous stand-up          & \texttt{08\_conclusions.tex} & Heading only (commented) \\
Single-wheel variant         & \texttt{08\_conclusions.tex} & Heading + comment \\
\bottomrule
\end{tabular}
\end{table}

\section{Comment-only stubs}

Down from twenty at the previous scan to one: the seven electrical
schematic-sheet placeholders (\texttt{11\_electrical\_detail.tex}) are now
live \texttt{\textbackslash phfig} figures with captions and per-pin tables
(A14), and the former Section~3 stubs no longer exist as a section.

\begin{table}[h]
\centering
\small
\begin{tabular}{L{4.6cm} L{9.0cm}}
\toprule
Location & Content \\
\midrule
\texttt{05\_cad\_design.tex}, \texttt{app:cad\_3d} in
\texttt{10\_cad\_detail.tex} & 3D cube drawings: exploded view; section view
through two orthogonal wheel modules. Drawing set not yet released --- waits
on A4. \\
\bottomrule
\end{tabular}
\end{table}

\section{Missing figures}

\begin{itemize}
  \item \texttt{04\_methodology.tex}, \texttt{sec:2dmodel} --- free-body
        diagram of the 2D pendulum, labelling $\theta_b$, $\theta_w$, $l$,
        $l_b$ and the pivot. Unchanged from the previous scan; this is still
        the only figure suggested but not present in the dynamics section.
  \item \texttt{05\_cad\_design.tex}, \texttt{app:cad\_3d} --- cube exploded
        view and the section view named in the comment-only stub above.
\end{itemize}

The seven electrical schematic sheets are no longer in this list: they render
as placeholder-art figures (Section~\ref{sec:tbd}, A14) rather than as
comments, so they are tracked as artwork-pending rather than as missing
figures.

\section{Defects found alongside the TBDs}

Historical record, retained for audit. All rows below predate the
restructure; renumbered section/table references are current.

\begin{table}[h]
\centering
\small
\begin{tabular}{c L{4.2cm} L{8.4cm}}
\toprule
\# & Location & Defect \\
\midrule
1 & Power architecture vs.\ BOM & \closed{} Battery capacity contradiction resolved against the procured pack: Tattu R-Line V5.0, \qty{22.2}{\volt}, \qty{1550}{\milli\ampere\hour}, \num{150}C, XT60. \\
2 & \texttt{tab:rw\_params} & \closed{} Missing \texttt{\textbackslash caption} added; now appears in the List of Tables. \\
3 & G6 hard-kill date vs.\ \texttt{tab:gates} & \closed{} Both now state 17~August. \\
4 & \texttt{02\_project\_organization\_tasks} & Task B4 says ``Raspberry Pi software architecture''; the rest of the document specifies Teensy~4.1 throughout. \\
5 & BOM vs.\ CAD part list & Brake servo quantity: BOM says 1--3, the CAD part list says 3 brackets. \\
6 & Author list vs.\ role/task tables & \closed{} Name spelling unified to ``Niccolo'' throughout. \\
7 & \texttt{02\_project\_organization\_tasks} & Task IDs C11--C14 appear under the ``3D Design'' block, after D1--D7. \\
8 & All section files & British/American spelling mixed: realise/realize, characterise/characterize, behaviour/behavior. \\
9 & \closed{} Hard-coded \texttt{\textbackslash newpage} & Removed with the section it sat in during the restructure. \\
10 & All tables & \closed{} Every float now uses \texttt{[htbp]}; no \texttt{[h]} float remains in \texttt{sections/}. \\
11 & \texttt{sections/} & \texttt{code.cpp} and \texttt{code.exe} are untracked in the repo; a compiled binary should not be committed. \\
\bottomrule
\end{tabular}
\end{table}

\section{Suggested order of work}

This mirrors, and should be kept in sync with, the Immediate Next Steps
subsection of \texttt{08\_conclusions.tex} (\texttt{sec:next\_steps}).

\begin{enumerate}
  \item \textbf{Close the packaging study} (A4) --- still the head of the
        register. $V_{\min}$ and the clearances release $L$, then $M$. A1, A2
        and A8 are populated but \emph{provisional}: they rest on a
        \qty{610}{\gram} structural allowance, so closing A4 moves them again
        as the next turn of the fixed-point iteration, not a re-opening of
        settled items.
  \item \textbf{Build and run the Simulink/Simscape environment}
        (\texttt{sec:sim\_environment}) against both prototypes' CAD, and
        re-derive the gains of A15 once A9 (plant ID) reports measured
        parameters.
  \item \textbf{Do the \ready{} items}, which need no upstream data: A3 (rail
        budget), A5 (fabrication routes), A6 (print settings), A11
        (connectors), A12 (wire gauges).
  \item \textbf{Decide the board layout} (\texttt{sec:layout}) to close the
        physical pin numbers of A14.
  \item \textbf{Fix the remaining defects} (Section~\ref{sec:tbd} above),
        which cost credibility out of proportion to the effort of fixing
        them.
  \item \textbf{The remaining \waiting{} items} (A7, A9, A10, A13, A15 and the
        three prose TBDs) close on hardware and cannot be accelerated on
        paper.
\end{enumerate}

\end{document}
